\chapter{Production Planning: From Demand to a Feasible Production Plan}
\label{ch:production-planning}

\chapterguide
  {You should know that Sales generates demand information and Supply Chain must turn that demand into physical products.}
  {explain the production-planning hierarchy, distinguish common production approaches, connect forecasts to capacity and inventory decisions, and explain the role of a bill of materials and manufacturing order.}

\begin{sourcebasis}
This chapter follows Tutorial 3 (Production Planning) and Chapter 4 of the
textbook. The tutorial asks students to select a production-planning approach,
calculate a sales forecast, and prepare an S\&OP plan. The Odoo build guide
provides the Road Bike bill of materials and manufacturing workflow.
\end{sourcebasis}

\section{What production planning is trying to decide}

Production planning coordinates demand, inventory, materials, capacity, and
time. The planner works down a hierarchy in which each level fixes a decision
the level below must respect, and each level covers a shorter horizon in
finer detail.

\begin{erpfig}[The planning hierarchy]{The planning hierarchy read as a pyramid.
Moving down, the horizon shortens and the detail sharpens: a yearly forecast for
a product family becomes a daily instruction for one
product.}{fig:planning-hierarchy}
\begin{tikzpicture}[x=1cm,y=1cm]
\foreach \i/\lvl/\q/\h/\d in {%
  0/{Sales forecast}/{What will customers want?}/{12 months}/{product family},
  1/{Sales \& Operations Planning}/{How much can we make, and when?}/{monthly}/{product family},
  2/{Demand management}/{How much of each product, each week?}/{weekly}/{single product},
  3/{Detailed scheduling}/{Which product runs on which day?}/{daily}/{one run}}{
  \pgfmathsetmacro{\y}{-\i*1.12}
  \pgfmathsetmacro{\w}{10.6-\i*0.9}
  \pgfmathsetmacro{\tw}{\w-0.55}
  \node[erpstep,minimum width=\w cm,minimum height=0.9cm,text width=\tw cm] (p\i) at (6.0,\y)
    {\lvl\ --- \emph{\q}};
  \node[erpchip,anchor=west] at (11.6,\y) {\h};
  \node[erpnote,anchor=west,text width=2.3cm] at (13.5,\y) {\d};}
\foreach \i in {0,...,2}{\pgfmathtruncatemacro{\j}{\i+1}\draw[erpflow] (p\i)--(p\j);}
\node[erpgood,minimum width=3.3cm,minimum height=0.8cm] (prod) at (3.5,-4.6) {Production};
\node[erpgood,minimum width=3.3cm,minimum height=0.8cm] (mrp)  at (8.5,-4.6) {MRP $\rightarrow$ Purchasing};
\draw[erpflow] (p3.south) -- ++(0,-0.28) -| (prod.north);
\draw[erpflow] (p3.south) -- ++(0,-0.28) -| (mrp.north);
\node[erpcap,anchor=west] at (11.6,0.78) {\bfseries horizon};
\node[erpcap,anchor=west] at (13.5,0.78) {\bfseries detail};
\end{tikzpicture}
\end{erpfig}

\section{Production approaches}

The most useful distinction is \emph{when} production happens relative to the
customer order. The point where the customer's order first touches the process
is called the decoupling point: everything upstream of it is driven by forecast,
everything downstream by a real order.

\begin{erpfig}[Where the customer order enters]{The same four activities, three
strategies. The orange marker is the decoupling point: work to its left is driven
by forecast, work to its right by a real order. Moving it left cuts inventory
risk and lengthens the customer's wait.}{fig:decoupling-point}
\begin{tikzpicture}[x=1cm,y=1cm]
\foreach \x/\t in {2.0/Design, 5.0/{Buy parts}, 8.0/Manufacture, 11.0/{Ship to customer}}{
  \node[erpcap] at (\x,3.05) {\bfseries \t};}
\foreach \i/\nm/\k in {0/{Make-to-stock}/4, 1/{Make-to-order}/3, 2/{Engineer-to-order}/1}{
  \pgfmathsetmacro{\y}{2.2-\i*1.55}
  \node[erpcap,anchor=east,align=right,text width=1.95cm] at (0.35,\y) {\bfseries \nm};
  \foreach \s in {1,2,3,4}{
    \pgfmathsetmacro{\x}{-1.0+\s*3.0}
    \ifnum\s<\k
      \node[erpquiet,minimum width=2.6cm,minimum height=0.72cm] at (\x,\y) {forecast};
    \else
      \node[erpgood,minimum width=2.6cm,minimum height=0.72cm] at (\x,\y) {customer order};
    \fi}
  \pgfmathsetmacro{\dx}{-1.0+\k*3.0-1.5}
  \draw[draw=UTPOrange,line width=1.3pt] (\dx,\y-0.52) -- (\dx,\y+0.52);
  \node[erpchipwarn,anchor=north] at (\dx,\y-0.6) {order enters here};}
\node[erpnote,anchor=west,text width=12.6cm,align=left] at (-0.3,-2.55)
  {\emph{Left} of the marker the company is guessing and carrying inventory risk.
   \emph{Right} of it the customer is waiting. West End Bicycles uses
   make-to-stock: Road Bikes are built before the customer order is delivered.};
\end{tikzpicture}
\end{erpfig}

\begin{erptable}
\begin{tabularx}{\linewidth}{@{}>{\bfseries\leavevmode\color{ERPNavy}}p{0.21\linewidth}p{0.31\linewidth}X@{}}
\toprule
\rowcolor{ERPPaleBlue!92}
Approach & Core idea & Typical trade-off \\
\midrule
Make-to-stock & Produce finished goods before a specific customer order, based on expected demand. & Fast fulfilment, but inventory risk if demand is wrong. \\
Make-to-order & Begin production after a customer order is received. & Lower finished-goods inventory, but longer customer lead time. \\
Engineer/configure-to-order & Design or configure significant product details after the order. & High customisation, but greater planning complexity and longer lead time. \\
\bottomrule
\end{tabularx}
\end{erptable}

\section{Sales forecasting}

A \term{sales forecast} estimates future demand. It is neither a promise nor a
confirmed order. Forecasts are needed because supply decisions require lead time.

The textbook's approach starts from last year's demand, removes the unusual
effect of last year's promotion, applies expected underlying growth, then adds
this year's planned promotion.

\begin{erpfig}[Building one month's forecast]{The four-step adjustment for a
single month. Skipping the first step would treat a one-off promotion as
permanent base demand and inflate every later
number.}{fig:forecast-build}
\begin{tikzpicture}[x=1cm,y=1cm]
\foreach \i/\lbl/\val/\c in {%
  0/{Last year's sales}/{5,500}/UTPMuted,
  1/{$-$ last year's promotion}/{$-$ 200}/UTPRed,
  2/{$=$ base demand}/{5,300}/UTPBlue,
  3/{$\times\,(1+g)$, $g=3\%$}/{$+$ 159}/UTPGreen,
  4/{$=$ base projection}/{5,459}/UTPBlue,
  5/{$+$ this year's promotion}/{$+$ 0}/UTPGreen,
  6/{$=$ forecast}/{5,459}/UTPGreen}{
  \pgfmathsetmacro{\y}{-\i*0.78}
  \node[erpdoc,anchor=west,minimum width=7.4cm,text width=7.0cm,align=left,
        minimum height=0.66cm,draw=\c!65] (b\i) at (1.6,\y) {\lbl};
  \node[erpchip,anchor=west,text width=1.5cm,align=center] at (9.3,\y) {\val};}
\node[erpcap,anchor=west,text width=1.5cm,align=left] at (0,-0.4) {\bfseries remove\\ the one-off};
\node[erpcap,anchor=west,text width=1.5cm,align=left] at (0,-2.75) {\bfseries grow the\\ base};
\node[erpcap,anchor=west,text width=1.5cm,align=left] at (0,-4.5) {\bfseries add the\\ new plan};
\node[erpnote,anchor=west,text width=4.0cm] at (11.2,-2.35)
  {June for Footloose Sport Shoes: no promotion is planned this June, so the
   forecast equals the grown base.};
\end{tikzpicture}
\end{erpfig}

For month $t$, with expected growth rate $g$:

\[
\text{Base}_t = \text{Prev.\ sales}_t - \text{Prev.\ promo}_t,
\qquad
\text{Forecast}_t = \text{Base}_t(1+g) + \text{New promo}_t.
\]

\begin{keyidea}
Forecasting must separate the underlying demand pattern from one-off causes.
Multiplying a promotion-inflated total by a growth rate compounds the error every
year.
\end{keyidea}

\section{Sales and Operations Planning}

\term{Sales and Operations Planning (S\&OP)} turns the forecast into a feasible
aggregate production plan. Inventory is the buffer that lets production differ
from demand in any single month.

\begin{erpfig}[The inventory balance]{S\&OP as a tank. Production fills it,
demand drains it, and the level at the end of the month is the level at the
start of the next. Capacity limits how fast the inlet can
run.}{fig:sop-balance}
\begin{tikzpicture}[x=1cm,y=1cm]
\draw[draw=UTPBlue,fill=UTPPaleBlue!55,rounded corners=1.5mm] (5.0,-0.2) rectangle (9.6,2.4);
\draw[draw=UTPBlue!35,fill=UTPPaleBlue] (5.05,-0.15) rectangle (9.55,1.15);
\node[erphead] at (7.3,0.5) {Inventory};
\node[erpnote] at (7.3,1.75) {level at month end};
\node[erpstep,minimum width=3.0cm,minimum height=0.9cm] (prod) at (1.4,2.4)
  {Production $P_t$\\\tiny\mdseries what we choose to make};
\node[erpgood,minimum width=3.0cm,minimum height=0.9cm] (dem) at (13.2,2.4)
  {Demand $F_t$\\\tiny\mdseries what customers take};
\draw[erpflow] (prod.east) -- (5.0,2.4) -- (5.0,1.6);
\draw[erpflow] (9.6,1.6) -- (9.6,2.4) -- (dem.west);
\node[erpwarn,minimum width=3.0cm,minimum height=0.8cm] (cap) at (1.4,0.6)
  {Capacity $C_t$\\\tiny\mdseries the ceiling on $P_t$};
\draw[erpslow] (cap.north) -- (prod.south);
\node[erpstore,minimum width=6.6cm,minimum height=0.85cm,draw=UTPNavy,fill=UTPNavy!7]
  at (7.3,-1.35) {$I_t = I_{t-1} + P_t - F_t$};
\node[erpnote,anchor=west,text width=3.6cm] at (10.6,0.5)
  {a positive balance is the buffer carried into next month};
\end{tikzpicture}
\end{erpfig}

A feasible plan must respect capacity:

\[
C_t = r \times h \times D_t,
\qquad
U_t = \frac{P_t}{C_t}\times 100\%,
\]

where $r$ is the production rate, $h$ the hours per day, and $D_t$ the working
days in the month. The plan therefore balances three competing pressures:
customer demand, inventory holding, and capacity.

\section{Demand management and detailed scheduling}

S\&OP is deliberately aggregate. \term{Demand management} breaks the aggregate
plan into individual products and shorter periods; \term{detailed scheduling}
then fixes the actual runs. Longer runs reduce changeovers but raise inventory;
shorter runs cut inventory but spend capacity on setup.

\begin{intuition}
S\&OP decides ``how much should the factory produce this month?'' Detailed
scheduling decides ``which product runs on which day, in what sequence?''
\end{intuition}

\section{The bill of materials connects planning to materials}

A \term{bill of materials (BOM)} is the product recipe: the components and
quantities needed to make one finished unit. It is what converts a production
quantity into a purchasing requirement.

\begin{erpfig}[The Road Bike bill of materials]{The supplied BOM, drawn as the
tree it actually is. The right-hand column shows the same structure multiplied
out for the manufacturing order of 5 bikes used throughout the
practicals.}{fig:road-bike-bom}
\begin{tikzpicture}[x=1cm,y=1cm]
\node[erpstep,minimum width=3.4cm,minimum height=0.95cm,draw=UTPGreen,fill=UTPPaleTeal]
  (top) at (3.2,2.6) {Road Bike\\\tiny\mdseries 1 finished unit};
\foreach \i/\c/\q in {0/{Aluminium Frame}/1, 1/{Rubber Tires}/2, 2/{Brake Set}/1,
                      3/{Gear Set}/1, 4/{Bicycle Chain}/1}{
  \pgfmathsetmacro{\y}{1.0-\i*0.82}
  \node[erpdoc,anchor=west,minimum width=3.4cm,text width=3.0cm,align=left,
        minimum height=0.68cm] (c\i) at (2.6,\y) {\c};
  \node[erpchip,anchor=west] at (6.2,\y) {$\times\,\q$};
  \draw[erplink] (1.5,2.15) |- (c\i.west);}
\draw[erplink] (1.5,2.15) -- (1.5,-2.28);
\node[erpcap,anchor=west,color=UTPMuted] at (0.2,-0.1) {\bfseries per unit};
% multiplied out
\draw[draw=UTPBlue!35,fill=UTPPaleBlue!40,rounded corners=2.5mm] (8.4,-2.72) rectangle (14.6,3.25);
\node[erpcap,color=UTPBlue] at (11.5,2.95) {\bfseries for a Manufacturing Order of 5};
\node[erpstep,minimum width=4.2cm,minimum height=0.8cm,draw=UTPGreen,fill=UTPPaleTeal]
  at (11.5,2.2) {5 Road Bikes produced};
\foreach \i/\c/\q in {0/{Aluminium Frame}/5, 1/{Rubber Tires}/10, 2/{Brake Set}/5,
                      3/{Gear Set}/5, 4/{Bicycle Chain}/5}{
  \pgfmathsetmacro{\y}{1.0-\i*0.82}
  \node[erpdoc,anchor=west,minimum width=3.2cm,text width=2.8cm,align=left,
        minimum height=0.68cm] at (8.8,\y) {\c};
  \node[erpchipwarn,anchor=west] at (12.2,\y) {$-\,\q$ from stock};}
\node[erpnote,anchor=north,text width=12.4cm] at (7.4,-3.05)
  {Completing the Manufacturing Order consumes the components on the left and
   creates the finished bikes on the right --- one event, two stock effects.};
\end{tikzpicture}
\end{erpfig}

MRP later extends this multiplication by considering inventory, scheduled
receipts, lot sizes, and lead times.

\section{Manufacturing orders execute the plan}

A plan says what should happen. A \term{manufacturing order (MO)} is the
instruction to produce a specific quantity. In an integrated system, completing
the MO is a stock event, not just a status change.

\begin{erpfig}[What completing an MO does to stock]{Before and after the
Manufacturing Order for 5 Road Bikes. This before/after pair is the evidence the
labs ask you to capture, and the quickest way to prove the manufacturing event
reached the shared inventory data.}{fig:mo-stock-change}
\begin{tikzpicture}
\begin{axis}[
  width=13.6cm, height=5.4cm,
  ybar=3pt, bar width=13pt,
  ymin=0, ymax=32,
  ylabel={units on hand},
  symbolic x coords={Frame,Tires,Brake Set,Gear Set,Chain,Road Bike},
  xtick=data, x tick label style={font=\scriptsize\color{UTPMuted}},
  enlarge x limits=0.11,
  legend style={at={(0.5,1.17)},anchor=north,legend columns=2},
  ymajorgrids, grid style={draw=UTPMuted!22},
  nodes near coords, every node near coord/.append style={font=\tiny\color{UTPMuted}},
]
\addplot[draw=UTPBlue,fill=UTPPaleBlue] coordinates
  {(Frame,10) (Tires,20) (Brake Set,10) (Gear Set,10) (Chain,10) (Road Bike,0)};
\addplot[draw=UTPGreen,fill=UTPPaleTeal] coordinates
  {(Frame,5) (Tires,10) (Brake Set,5) (Gear Set,5) (Chain,5) (Road Bike,5)};
\legend{before the MO, after the MO is marked Done}
\end{axis}
\end{tikzpicture}
\end{erpfig}

\begin{odoobox}
In the West End Bicycles guide: create a Road Bike BOM for one unit; create a
Manufacturing Order for 5; assign the Production Planner as responsible; confirm
and check component availability; produce and mark the MO Done; then verify that
component stock fell and Road Bike on-hand stock rose to 5.
\end{odoobox}

\section{Safety stock and the cost of uncertainty}

\term{Safety stock} is extra inventory held to reduce the risk of a stockout
when demand or supply is uncertain. More safety stock buys service protection
and costs holding expense and tied-up cash --- and the exchange rate between the
two is not linear.

\begin{erpfig}[Diminishing returns on safety stock]{Service level against safety
stock for a representative demand distribution. The first units of buffer are
cheap protection; the last few percent of service are very expensive, which is
why 100\% service is almost never the target.}{fig:safety-stock-curve}
\begin{tikzpicture}
\begin{axis}[
  width=12.8cm, height=5.6cm,
  xlabel={safety stock held (units)}, ylabel={service level (\%)},
  xmin=0, xmax=100, ymin=45, ymax=102,
  ymajorgrids, grid style={draw=UTPMuted!22},
  legend style={at={(0.99,0.03)},anchor=south east},
]
\addplot[draw=UTPBlue,smooth,domain=0:100,samples=80]
  {50 + 50*(1-exp(-x/26))};
\addlegendentry{service level achieved}
\addplot[draw=UTPOrange,dashed,domain=0:100,samples=2] {95};
\addlegendentry{95\% target}
\draw[draw=UTPGreen,line width=1pt,dashed] (axis cs:78,45) -- (axis cs:78,95);
\node[erpchip,anchor=west,font=\tiny] at (axis cs:16,72) {cheap protection};
\node[erpchipwarn,anchor=east,font=\tiny] at (axis cs:76,62) {expensive last few \%};
\end{axis}
\end{tikzpicture}
\end{erpfig}

Forecast accuracy matters because a worse forecast needs a bigger buffer for the
same service level. ERP does not remove uncertainty; it improves the information
used to manage it.

\begin{pitfall}
High capacity utilisation is not the same as good planning. A factory can keep
machines busy producing goods nobody ordered. The objective is a sensible balance
between demand, capacity, setup cost, and inventory cost.
\end{pitfall}

\begin{tryit}
Choose one of the ventures suggested in Tutorial 3 --- furniture, cookies,
printing, or your own. Decide whether make-to-stock, make-to-order, or a more
customised approach fits, then explain what would go wrong under one of the
others. Mark where the decoupling point sits in each case.
\end{tryit}

\begin{quickcheck}
\begin{enumerate}
  \item Why can production differ from the sales forecast in a given month?
  \item What information does a BOM provide that a sales forecast does not?
  \item What stock changes should occur when a Manufacturing Order for 5 Road Bikes is completed?
\end{enumerate}
\end{quickcheck}

\begin{chapterrecap}
\begin{itemize}
  \item Planning descends a hierarchy from yearly product families to daily runs (\figref{fig:planning-hierarchy}).
  \item The production approach is really a choice of where the customer order enters the process (\figref{fig:decoupling-point}).
  \item Forecasting removes one-off promotions before applying growth (\figref{fig:forecast-build}).
  \item S\&OP balances production, demand, and inventory subject to capacity (\figref{fig:sop-balance}).
  \item A BOM converts a production quantity into a component requirement; the MO makes it a stock event (\figref{fig:road-bike-bom} and \figref{fig:mo-stock-change}).
\end{itemize}
\end{chapterrecap}
